% Part: history
% Chapter: biographies
% Section: rozsa-peter

\documentclass[../../../include/open-logic-section]{subfiles}

\begin{document}

\olfileid{his}{bio}{pet} 

\olsection{रोझा पेटर}

\olphoto{peter-rozsa}{रोझा पेटर}

रोझा पेटर यांचा जन्म १७ फेब्रुवारी १९०५ रोजी हंगेरीतील बुडापेस्ट येथे रोझा पोलित्झर या नावाने झाला. पुनरावर्ती फलनांवरील कार्यासाठी त्या विशेष ओळखल्या जातात; पुनरावर्तन उपपत्ती हे क्षेत्र निर्माण होण्यात हे कार्य अत्यावश्यक ठरले.

पेटर यांचे बालपण राजकीय उलथापालथीच्या काळात गेले; त्या किशोरवयीन असताना पहिले महायुद्ध सुरू होते. तरीही त्यांना बुडापेस्टमधील प्रतिष्ठित मारिया तेरेझिया मुलींच्या शाळेत शिकता आले आणि १९२२ मध्ये त्यांनी तेथील शिक्षण पूर्ण केले. त्यानंतर त्यांनी बुडापेस्टच्या पाझमान्य पेटर विद्यापीठात शिक्षण घेतले; पुढे त्याचे नाव लोरांड एटव्होश विद्यापीठ झाले. वडिलांच्या आग्रहामुळे त्यांनी सुरुवातीला रसायनशास्त्र घेतले, पण नंतर गणिताकडे वळल्या आणि १९२७ मध्ये पदवीधर झाल्या. माध्यमिक शाळेत गणित शिकवण्याची पात्रता असूनही त्या काळची आर्थिक परिस्थिती भीषण होती; महामंदीने जागतिक अर्थव्यवस्था प्रभावित झाली होती. म्हणून त्यांनी शिकवण्या आणि खाजगी गणित अध्यापन अशी प्रसंगोपात्त कामे केली. अखेरीस गणितातील पदव्युत्तर शिक्षणासाठी त्या पुन्हा विद्यापीठात आल्या. सुरुवातीला त्यांचा संख्या उपपत्तीवर काम करण्याचा विचार होता; पण आपले निष्कर्ष आधीच सिद्ध झाले आहेत हे कळल्यावर त्यांनी गणितच सोडून देण्याचा जवळजवळ निर्णय घेतला. त्यांना गोडेल यांच्या अपूर्णता प्रमेयांवर काम करण्यास प्रोत्साहन मिळाले आणि नकळत त्यांनी त्यांचे काही निष्कर्ष वेगळ्या पद्धतीने सिद्ध केले. त्यामुळे त्यांचा आत्मविश्वास परत आला. डेव्हिड हिल्बर्ट यांच्या पायाभूत कार्यक्रमातून प्रेरणा घेऊन पेटर यांनी पुनरावर्तन उपपत्तीवरील आपले पहिले लेख लिहिले. १९३५ मध्ये त्यांनी डॉक्टरेट मिळवली आणि १९३७ मध्ये त्या \emph{Journal of
  Symbolic Logic} च्या संपादक झाल्या.

पेटर यांचे आरंभीचे लेख पुनरावर्ती फलनांच्या उपपत्तीचे संस्थापक योगदान म्हणून व्यापकपणे मानले जातात. \cite{Peter1935a} मध्ये त्यांनी पुनरावर्तनाच्या विविध प्रकारांमधील संबंध तपासले. \cite{Peter1935b} मध्ये त्यांनी पुनरावर्ती पद्धतीने परिभाषित केलेले एक विशिष्ट फलन आदिम पुनरावर्ती नाही, हे दाखवले. यामुळे विल्हेल्म ॲकरमन यांचा आधीचा निष्कर्ष सोपा झाला. पेटर यांचे हे सुलभ केलेले फलन आज अनेकदा ॲकरमन फलन म्हणतात; काही वेळा अधिक अचूकपणे ॲकरमन–पेटर फलनही म्हणतात. पुनरावर्ती फलनांच्या उपपत्तीवरील पहिले पुस्तक त्यांनी लिहिले \citep{Peter1951}.

त्यांच्या कार्याचे महत्त्व आणि प्रभाव मोठा असूनही पेटर यांना १९४५ पर्यंत पूर्णवेळ अध्यापनाचे पद मिळाले नाही. दुसऱ्या महायुद्धात नाझींनी हंगेरीवर कब्जा केला असताना यहुदीविरोधी कायद्यांमुळे त्यांना शिकवण्याची परवानगी नव्हती. १९४४ मध्ये सरकारने बुडापेस्टमध्ये ज्यू लोकांसाठी बंदिस्त वस्ती निर्माण केली; ती शहराच्या इतर भागांपासून अलग ठेवली गेली होती आणि सशस्त्र रक्षक तिच्यावर पहारा देत होते. १९४५ मध्ये ही वस्ती मुक्त होईपर्यंत पेटर यांना तेथे राहण्यास भाग पाडले गेले. त्यानंतर त्यांनी बुडापेस्ट टीचर्स ट्रेनिंग कॉलेजमध्ये आणि १९५५ पासून एटव्होश लोरांड विद्यापीठात अध्यापन केले. ‘गणितातील अकादमिक डॉक्टर’ हा दर्जा मिळवणाऱ्या त्या पहिल्या हंगेरियन महिला गणितज्ञ आणि हंगेरियन विज्ञान अकादमीवर निवडून जाणाऱ्या पहिल्या महिला होत्या.

पेटर या गणित उत्कटतेने शिकवणाऱ्या अध्यापिका म्हणून ओळखल्या जात. केवळ व्याख्यान देण्याऐवजी विद्यार्थ्यांसोबत गणिती समस्यांचे स्वरूप आणि सौंदर्य शोधणे त्यांना अधिक आवडत असे. म्हणून विद्यार्थी त्यांना प्रेमाने “रोझा मावशी” म्हणत. १९७७ मध्ये वयाच्या ७१व्या वर्षी त्यांचे निधन झाले.

\begin{reading}
अधिक चरित्रात्मक वाचनासाठी \citep{Oconnor2014} आणि \citep{Andrasfai1986} पाहा. पेटर यांची एक छोटी मुलाखत \citet{Tamassy1994} यांनी घेतली होती. गणिताविषयी रंजक वाचनासाठी पेटर यांचे \emph{Playing With Infinity} हे पुस्तक \citep{Peter2010} पाहा.
\end{reading}

\end{document}
